\chapter{Aplicación a Datos Reales}
\label{04_00_00_aplicacion}

% Esqueleto (2026-10-12). Misma estructura que los resultados de simulación (§03_08); las corridas son
% 101 (rezago 1) y 102 (rezagos {1, 7}) de notebooks/reales/, con las mismas salidas que las 30x.

\section{Datos}
\label{04_01_00_datos}
% Tráfico I-94 (Metro Interstate Traffic Volume): curvas diarias, ventana 2015-07-01..2018-09-30, 70/30.
% Representación B-spline 10/3, FPCA en train, M = 10. Fuente: 101_01 §2–§3.

\section{Diseño}
\label{04_02_00_diseno}
% Lo que difiere del estudio de simulación: dos especificaciones (rezago 1; rezagos {1, 7} = día anterior y
% mismo día de la semana anterior), FAR(7) completo, RF con subconjuntos {1} y {1, 7}, atau = 2, 1 cadena,
% ventana de detalle por calendario (10 días de prueba con más no laborables). Cuadro con la configuración.

\section{Resultados}
\label{04_03_00_resultados}

\subsection{Error puntual}
\label{04_03_01_error_puntual}
% Cuadro: MAE de los tres predictores, FAR y RF para las dos especificaciones (M = 10). Fuente: 98_* de 101_05 y 102_05.

\subsection{Intervalos de predicción}
\label{04_03_02_intervalos}
% Cuadro: Winkler, MPIW, PICP de las tres construcciones y del FAR; % de ventanas que gana al FAR. Fuente: 99_*, 111_*.

\subsection{Comportamientos predictivos}
\label{04_03_03_comportamientos}
% Distribución de K en test; K contra día de la semana y feriados (65_*, 104_*); ventana de detalle (112_*).

\subsection{Pronóstico por comportamiento}
\label{04_03_04_pronostico}
% Cuadro del _06 (119_*): pronóstico, cota, esperanza, FAR, RF; todos / K >= 2; acierto por tramos de p_1.

\section{Discusión}
\label{04_04_00_discusion}
